% titel: Blatt 7, A5: Partialdrücke (He, Methan, Ethan) und Gesamtdruck
% Teil II: der durchgerechnete Weg.

\subsection*{Gegeben und gesucht}
$V=40\ \mathrm{L}$, $T=300\ \mathrm{K}$; $m_{\mathrm{He}}=1.0$ g,
$m_{\mathrm{CH_4}}=5.0$ g, $m_{\mathrm{C_2H_6}}=4.0$ g. Gesucht:
Partialdr\"ucke der drei Gase und Gesamtdruck.

\subsection*{Ansatz}
F\"ur jedes Gas $n_i=m_i/M_i$ ausrechnen, dann $p_i=n_i RT/V$; Summe
$=$ Gesamtdruck.

\subsection*{Rechnung}
Stoffmengen ($M_{\mathrm{He}}=4.00$, $M_{\mathrm{CH_4}}=16.04$,
$M_{\mathrm{C_2H_6}}=30.07$ in $\mathrm{g\,mol^{-1}}$):
\begin{align*}
  n_{\mathrm{He}} &= \frac{1.0}{4.00} = 0.250\ \mathrm{mol},\\
  n_{\mathrm{CH_4}} &= \frac{5.0}{16.04} = 0.3117\ \mathrm{mol},\\
  n_{\mathrm{C_2H_6}} &= \frac{4.0}{30.07} = 0.1330\ \mathrm{mol}.
\end{align*}
Partialdr\"ucke ($R=0.0821\ \mathrm{L\,atm\,mol^{-1}K^{-1}}$):
\begin{align*}
  p_{\mathrm{He}} &= 0.250\cdot\frac{0.0821\cdot 300}{40}
                    = 0.154\ \mathrm{atm},\\
  p_{\mathrm{CH_4}} &= 0.3117\cdot\frac{0.0821\cdot 300}{40}
                    = 0.192\ \mathrm{atm},\\
  p_{\mathrm{C_2H_6}} &= 0.1330\cdot\frac{0.0821\cdot 300}{40}
                    = 0.0819\ \mathrm{atm}.
\end{align*}
Gesamtdruck:
\[
  p_{\text{ges}} = 0.154 + 0.192 + 0.0819 = 0.428\ \mathrm{atm}.
\]

\begin{ergebnis}
$p_{\mathrm{He}}=0.154$ atm, $p_{\mathrm{CH_4}}=0.192$ atm,
$p_{\mathrm{C_2H_6}}=0.0819$ atm, Gesamtdruck $0.428$ atm.
\end{ergebnis}

\begin{kontrolle}
Anteile: $x_{\mathrm{He}}=0.250/0.6947=0.360$,
$x_{\mathrm{CH_4}}=0.449$, $x_{\mathrm{C_2H_6}}=0.191$. Ein Stoffmengen-
und Druck-Anteil decken sich; Summe der Partialdr\"ucke $=$ Gesamtdruck.
Plausibel, da ideal (Druck $<1$ atm bei 300 K in 40 L).
\end{kontrolle}

\begin{merksatz}
Partialdruck $p_i = n_i RT/V$, Gesamtdruck $=$ Summe der
Partialdr\"ucke. Bei idealen Gasen ist der Partialdruck direkt proportional
zum Stoffmengenanteil.
\end{merksatz}
